\documentclass[11pt]{article}
\usepackage[letterpaper,margin=1in]{geometry}
\usepackage[T1]{fontenc}
\usepackage[utf8]{inputenc}
\usepackage{microtype}
\usepackage[hidelinks=false,colorlinks=true,urlcolor=blue]{hyperref}
\setlength{\parindent}{0pt}
\setlength{\parskip}{6pt}
\pagestyle{empty}

\begin{document}

[Date]

AI's 10 to Watch Selection Committee\\
IEEE Intelligent Systems

Dear Members of the Selection Committee,

I am pleased to strongly recommend Rahul Gupta for IEEE Intelligent Systems' AI's 10 to Watch. I have worked with Rahul through our joint work at Amazon over the course of several years, and I have seen first-hand his ability to balance industrial innovation, rigorous research, and support for the academic AI community.

Rahul is one of the uncommon industry scientists who can identify a responsible-AI problem in deployed systems, formulate it as a crisp scientific question, and then build the organizational mechanisms needed to make the answer useful. This combination is visible across his trajectory at Amazon: he founded Responsible AI science for Alexa, led Responsible AI for Amazon Nova, and now serves as Responsible AI lead for Amazon's AGI / Frontier AI Assets organization.

My view of Rahul is grounded in evidence from our collaborations. We have worked together on trustworthy NLP and responsible-AI research including FLIRT: Feedback Loop In-context Red Teaming; work on the steerability of large language models toward data-driven personas; research on slang and informal language processing in large language models; fairness, privacy, and model-evaluation projects; and broader efforts around foundation-model safety and responsible deployment. These collaborations showed me Rahul's originality as a research thinker. He is especially good at seeing where existing benchmarks or governance mechanisms are too shallow, and then proposing evaluations that better reflect real-world model behavior.

Rahul's research record is strong by conventional academic measures: 200+ papers in top venues, more than 5,000 citations, and contributions across ACM FAccT, ACL, EMNLP, NeurIPS, ICLR, ICASSP, and related venues. His work includes BOLD for bias measurement in open-ended language generation, fairness measurement for text classifiers, methods for fairness in large-scale natural-language-understanding systems, FLIRT for red-teaming language models, privacy-preserving NLP methods, and recent work on CBRN uplift evaluation and agentic benchmark design. These are not disconnected papers; together they form a coherent program around measuring and improving trustworthy AI systems.

Rahul's industrial impact is equally important. In Alexa, he helped move fairness and privacy research into deployed product mechanisms, improving experiences for underserved customer populations whose needs can be hidden by aggregate model performance. In Nova and AGI, he has led work on safety evaluation, red-teaming, privacy-sensitive behavior, mitigation, external assessment, and launch readiness for frontier foundation models. Amazon Science has publicly identified him as senior science manager and RAI lead for Amazon's AGI organization in its article on Amazon's responsible-AI pipeline. Third-party assessments of Amazon Nova Premier's safety provide further evidence that this kind of responsible-AI leadership has real product and public-trust consequences.

Rahul is also an inventor of 25+ patented or patent-pending technologies. His Alexa-era patents and Nova-era filings show how his ideas become mechanisms for production AI: cohort determination, model configuration, user-data processing, NLP component evaluation, generative language models, and synthetic-data generation. This matters for AI's 10 to Watch because Rahul's contributions are not only influential in papers; they also shape deployed AI systems used at scale.

Finally, Rahul has made a serious investment in academic community building. He has helped organize TrustNLP, Trustworthy Speech Processing, and the SemEval Unlearning Challenge. He also led the Amazon Nova AI Challenge, which attracted more than 100 university applications and supported 10 university teams working on trusted AI through substantial research sponsorships, AWS credits, mentoring, and prize opportunities. I have found Rahul particularly ingenious in creating research programs that are useful both to industry practitioners and to academic researchers.

Rahul's career is still early, but his direction is clear. He is helping define what responsible AI leadership should look like for frontier systems: technically deep, empirically rigorous, deployable, and connected to the broader research community. I believe he is exactly the kind of researcher IEEE Intelligent Systems intended to recognize through AI's 10 to Watch, and I recommend him strongly.

Sincerely,

Richard Zemel\\
Trianthe Dakolias Professor of Engineering and Applied Science; Professor of Computer Science, Columbia University\\
Email / signature placeholder

\end{document}
